\chapter{Enanas Blancas y Estrellas de Neutrones: Reseña Histórica}

\begin{itemize}
 \item \emph{1844, Friedrich Wilhem Bessel}: Usó la técnica de paralaje estelar, diseñada por él, para encontrar la distancia a Sirius. Pero también descubrió perturbaciones en su órbita, por lo que dedujo que Sirius es un sistema binario, con periodo 50 años.
\item \emph{1862, Alvan Graham Clark}: Con el telescopio más potente de su época, logró observar la estrella compañera de Sirius A, ahora llamada Sirius B (se encontraba en su apoastro). Determinó sus parámetros: $L_A=23.5L_{\odot}$ y  $L_B=0.03L_{\odot}$. $M_A=2.3M_{\odot}$ y $M_B =1.0M_{\odot}$.
\item \emph{1915, Walter Adams}: Mediante técnicas espectroscópicas, descubrió que Sirius B es una estrella blanca-azulada muy caliente, $T_B=27000\,K$, mientras que $T_A=9910\,K$. Usando la ley de Stefan-Boltzmann, $R_B=5.5\cdot10^{6}\approx0.008R_{\odot}$. !`Sirius B es una estrella con la masa del Sol confinada a un volumen más pequeño que la Tierra!.$\;\;\Rightarrow\;\;\rho_B=3.0\cdot10^{9}\;kg\cdot m^{-3}$! Los demás astrónomos los consideran resultados ``absurdos''.
\item \emph{1926, W.S. Adams}: Midió los redshifts gravitacionales de las líneas espectrales emitidas por Sirius B, determinando en forma independiente el radio de la estrella, sirviendo de comprobación de la naturaleza compacta de este objeto y también como un test para la teoría de relatividad general.
\item \emph{Agosto 1926, Dirac} formula la estadística de Fermi-Dirac.
\item \emph{Diciembre 1926, R.H. Fowler}: Usó la estadística de Fermi-Dirac para explicar las enanas blancas: la presión de degeneración de los electrones mantiene el equilibrio en contra de la gravedad.
\item \emph{1930, S. Chandrasekhar}: Calculó modelos de enanas blancas tomando en cuenta efectos relativistas en la ecuación de estado degenerada de los electrones, descubriendo que ninguna enana blanca puede ser más masiva que $\sim1.2 M_{\odot}$ (conocida como Masa de Chandrasekhar)
\item \emph{1932, L.D. Landau} Dió una explicación elemental del límite de Chandrasekhar.
\item \emph{1932, James Chadwick}: Descubre el neutrón.
\item \emph{1934, Waalter Baade y Fritz Zwicky}: Proponen la existencia de estrellas de neutrones. Identifican una nueva clase de objetos astronómicos denominados supernovas. Sugieren que una supernova puede ser creada por el colapso de una estrella normal a una estrella de neutrones.
\item \emph{1939, Oppenheimer y Volkoff}: Realizan los primeros cálculos detallados de la estructura de estrellas de neutrones usando la teoría general de la relatividad. Proponen la idea que los núcleos de neutrones en estrellas normales pueden ser fuente de energía estelar (la fusión termonuclear aún no era entendida).
\item \emph{1939, Oppenheimer y Snyder}: Calculan el colapso de una esfera homogénea de un fluido sin presión usando relatividad general, y descubren cómo la esfera pierde comunicación con el resto del universo. Es el primer cálculo de cómo un agujero negro se puede formar.
\item \emph{1942, Duyvendak, Mayall, Oort, Baade y Minkowskii}: Deducen que la nebulosa del Cangrejo es el remanente de la supernova observada por astrónomos chinos en 1054. Observan en el centro de ella una estrella que identifican como el remanente de la estrella que explotó en dicho año.
\item \emph{1949, Kaplan}: Deriva los efectos de la teoría general de la relatividad sobre las curvas masa-radio para enanas blancas masivas. Deduce que la teoría de relatividad general induce una inestabilidad para un radio menor de $\sim1.1\cdot10^{3}\;km$.
\item \emph{1958, Schatzman, Wakano, Harrison y Wheeler}: Incorporan el decaimiento beta inverso para la ecuación de estado para la materia de enanas blancas. Muestran que este efecto induce una inestabilidad dinámica para enanas blancas con masa de $\sim M_{\odot}$ y radio menor a $\sim4\cdot10^{3}\;km$.
\item \emph{1963, Hoyle and Fowler}: Proponen la idea de estrellas supermasivas, calcularon sus propiedades, y sugirieron que podrían estar asociadas con núcleos galácticos activos y quásars.
\item \emph{1963-1964, Chandrasekhar y Feynman}: Desarrollan la teoría general relativista de las pulsaciones estelares, y Feynman la usó para mostrar que estrellas supermasivas, aunque newtonianas en su estructura, están sujetas a una inestabilidad debida a Relatividad General.
\item \emph{1967, Hewish et al.}: Descubrimientos de los pulsares.
\item \emph{1968, Gold}: Propone que los púlsares son estrellas de neutrones rotantes.
\item \emph{1968-1969 varios}: Descubrimientos casi simultáneos de los pulsares del Cangrejo (la misma estrella identificada en 1942) y Vela en remanentes conocidos de supernovas. Proveen evidencia para la formación de estrellas de neutrones y pulsares en explosiones de supernovas.
\item \emph{1975, Hulse y Taylor}: Descubrimiento del primer pulsar binario. Miden la masa de una estrella de neutrones y, a la vez, proveen de un test para la teoría general de la relatividad.
\end{itemize}